\CourseSection{Design and Production Approach}\label{sec:design}
\CourseSubsection{Retroencabulator development}
The autonomous golf pushcart will be developed through system design, subsystem prototyping, integration, testing, and refinement. The team has already obtained some technical components and will attempt to purchase a used motorized golf cart as the base platform. The Raspberry Pi, sensors, motor-control electronics, and smart-bag system will be integrated onto the cart, with development focused on motor control, sensor-based golfer following and obstacle detection, RFID-based club identification and tracking, and shot-data collection. A controller will be used to process sensor inputs and control the cart's movement during golfer following and obstacle avoidance. Each subsystem will be tested independently before being integrated into the complete system. Testing will begin in a controlled environment before progressing to outdoor conditions involving different walking speeds, obstacles, and terrain, with results used to refine the system and prioritize the MVP requirements.

\CourseSubsection{Milestones and division of labor}
\Guidance{Plan at least five technical milestones. Give each milestone a date, an observable completion condition, and one responsible member as the primary lead. Course report and presentation deadlines do \textbf{not} belong here. This table tracks technical progress!}

\begin{singlespace}
\setlength{\LTcapwidth}{\linewidth}
\begin{longtable}{@{}L{0.09\linewidth} L{\dimexpr0.5500000\linewidth-6.0000000\tabcolsep\relax} L{0.18\linewidth} L{0.18\linewidth}@{}}
\caption{Technical milestones and ownership.}\label{tab:proposal-milestones}\\
\rowcolor{MidGray}
\TableHeader{ID} & \TableHeader{Milestone} & \TableHeader{Due date} & \TableHeader{Lead}\\
\midrule
\endfirsthead
\endhead
\bottomrule
\endfoot
% Add or replace data rows below.
M1 & Motorized control of the golf cart & October 30th, 2026 & Connor \\
\addlinespace[3pt]
M2 & Raspberry Pi integration and hardware intergration of the necessary sensors & November 20th, 2026 & Vince \\
\addlinespace[3pt]
M3 & RFID-Based Club Identification and Tracking & December 18th, 2026 & Jp \\
\addlinespace[3pt]
M4 &  Obstacle detection/avoidance/emergency stop functionality & January 22nd, 2027 & Vince \\
\addlinespace[3pt]
M5 & User Interface Design and login system for golfer profile & February 5th, 2027 & Jp \\
\addlinespace[3pt]
M6 &  Shot-Data Collection and Club Recommendation System & February 19th, 2027 & Connor \\
\addlinespace[3pt]
M7 & Full Subsystem Integration and System Refinement & March 1, 2027 & All members \\
\addlinespace[3pt]
M8 & Testing and validation of retroencabulator complete & March 15, 2027 & All members \\
\end{longtable}
\end{singlespace}

\Guidance{Add one or two explanatory sentences per milestone, and describe completion evidence and dependencies where applicable.}

\CourseSubsection{Contingencies}
\Guidance{Identify likely technical or scheduling obstacles, the trigger for taking corrective action, any alternative approaches, and the expected effect on time and cost.}
The main project risks include difficulty obtaining a suitable motorized cart, unreliable sensor-based golfer following, limitations in RFID detection, motor issues, and delays during system integration. If a component fails to meet requirements during testing, the team will modify the software or use an alternative sensor, controller, or RFID configuration. If RFID tags cannot be reliably detected within the golf bag, the team will test alternative identification tags such as NFC tags for example. If a suitable motorized golf cart cannot be obtained, a different cart or drive platform will be considered. Outdoor testing may be moved to a controlled environment if weather or terrain prevents reliable testing. If the project falls behind schedule, the team will prioritize the MVP features of basic cart movement, golfer following, obstacle detection, emergency stopping, and RFID-based club identification and tracking, while postponing more advanced features such as shot-based club recommendations or mobile-app functionality. These alternatives may add some cost or development time but will help ensure that the core system is completed.
\clearpage
